\clearpage
\fancyhead[L]{\small JUNO.}
\hypertarget{native-268-lower}{}
\begin{center}
JUNO.
\end{center}

\noindent\rule{\linewidth}{0.4pt}

\begin{center}
Monatliche Correspondenz. Vol. XI. pp. 475–482. May 1805.
\end{center}

\noindent\rule{\linewidth}{0.4pt}

Here we give a few additional observations of \emph{Juno} made at more distant observatories. They came to our attention only later and cover the final part of this period of visibility. Some could no longer be made on the meridian and were possible only away from it, with equatorial instruments. It was chiefly because of the bright moonlight in the first days of February that \emph{Juno} could no longer be seen and could not then be found again. Only Dr. \textsc{Gauss} was fortunate enough to obtain a further observation of Juno on 20 February. It is the last observation of Juno that anyone reported to us.

We also received from \textsc{Piazzi} in Palermo a series of meridian observations made with his whole circle. We give them below.
\SourceLeaf{269}
\begin{center}
\begingroup\fontsize{9}{11}\selectfont\setlength{\tabcolsep}{3pt}\renewcommand{\arraystretch}{1.1}
\begin{tabular}{r|r|r|r}
1805 & \shortstack{Mean time\\in Palermo} & \shortstack{Apparent right\\ascension of $\JunoSymbol$} & \shortstack{Apparent southern\\declination of $\JunoSymbol$}\\ \hline
22 Nov. & $7^{\mathrm u}\ 42^{\mathrm m}\ 39^{\mathrm s}9$ & $357^\circ\ 12'\ 28''5$ & $.\ .\ .\ .\ .$ \\
24 & $7\ 36\ 12.5$ & $357\ 33\ 36.6$ & $10^\circ\ 49'\ 34''9$ \\
25 & $7\ 33\ 1.5$ & $357\ 44\ 54.9$ & $10\ 47\ 21.5$ \\
26 & $7\ 29\ 52.3$ & $357\ 56\ 36.9$ & $10\ 44\ 49.9$ \\
28 & $7\ 23\ 37.0$ & $358\ 20\ 49.5$ & $10\ 38\ 59.5$ \\
29 & $7\ 20\ 31.4$ & $358\ 33\ 26.1$ & $10\ 35\ 43.2$ \\
4 Dec. & $7\ 5\ 28.2$ & $359\ 42\ 42.3$ & $10\ 15\ 52.4$ \\
10 & $6\ 47\ 21.8$ & $1\ 5\ 11.9$ & $.\ .\ .\ .\ .$ \\
\end{tabular}\endgroup
\end{center}

\par\indent There seems to be an error in the observation of 10 December, since it differs considerably from the position calculated from the elements. Dr. \textsc{Gauss} carefully compared these observations with his fourth set of elements and obtained the following results.

\begin{center}
\begingroup\fontsize{9}{11}\selectfont\setlength{\tabcolsep}{3pt}\renewcommand{\arraystretch}{1.1}
\begin{tabular}{r|r|r|r|r}
1805 & \shortstack{Calculated right\\ascension\\of $\JunoSymbol$} & \shortstack{Calculated southern\\declination\\of $\JunoSymbol$} & \multicolumn{2}{c}{\shortstack{Errors in Dr. \textsc{Gauss}’s\\fourth set of elements\\in R.A.\qquad in decl.}}\\ \hline
22 Nov. & $357^\circ\ 12'\ 11''6$ & $.\ .\ .\ .\ .$ & $-16''9$ & $.\ .\ .$ \\
24 & $357\ 33\ 28.3$ & $10^\circ\ 49'\ 47''3$ & $-8.3$ & $+12''4$ \\
25 & $357\ 44\ 42.4$ & $10\ 47\ 25.9$ & $-12.5$ & $+4.4$ \\
26 & $357\ 56\ 18.3$ & $10\ 44\ 55.0$ & $-18.6$ & $+5.1$ \\
28 & $358\ 20\ 39.6$ & $10\ 39\ 5.0$ & $-9.9$ & $+5.5$ \\
29 & $358\ 33\ 22.3$ & $10\ 35\ 55.8$ & $-3,8$ & $+12.6$ \\
4 Dec. & $359\ 42\ 27.6$ & $10\ 16\ 9.3$ & $-14.7$ & $+16.9$ \\
10 &  &  &  &  \\
\end{tabular}\endgroup
\end{center}

\Needspace{14\baselineskip}
\par\indent These differences are so small that they scarcely seem to call for an improvement of the fourth set of elements for Juno. All our older planetary tables show differences in right ascension that are just as large and perhaps often larger still. Even so, the diligent Dr. \textsc{Gauss} based the following improved fifth set of elements for Juno on these observations and on one he made himself on 20 February.

\begin{center}
\begingroup\fontsize{8}{10}\selectfont\setlength{\tabcolsep}{3pt}\renewcommand{\arraystretch}{1.1}
\begin{tabular}{lr}
Epoch 1805, Seeberg meridian . . . . & $42^\circ\ 32'\ 36''$ \\
\multicolumn{2}{l}{$\left.\begin{array}{lr}\text{\AscendingNode\ 1805 . . . . . . . . . . . . . .}&171\ 4\ 15.6\\\text{Aphelion . . . . . . . . . . . . .}&233\ 11\ 39\end{array}\right\}\begin{array}{l}\text{assumed stationary}\\\text{relative to the stars}\end{array}$} \\
Inclination of the orbit . . . . . . . . . . & $13\ 3\ 38$ \\
Daily tropical motion . . . . . . . & $815''9595$ \\
Annual . . . . . . . . . . . . . . & $82^\circ\ 43'\ 45''2$ \\
Eccentricity . . . . . . . . . . . . & $0.254236$ \\
Logarithm of the semiaxis . . . . . . & $0.4256078$ \\
\end{tabular}\endgroup
\end{center}

\par\indent Here is the observation by Dr. \textsc{Gauss} just mentioned.

\begin{center}
\begingroup\fontsize{9}{11}\selectfont\setlength{\tabcolsep}{3pt}\renewcommand{\arraystretch}{1.1}
\begin{tabular}{c|c|c}
Mean time & R.A. $\JunoSymbol$ & Decl. bor. $\JunoSymbol$\\ \hline
1805, 20 Feb. $7^{\mathrm u}\ 11^{\mathrm m}\ 12^{\mathrm s}$ & $30^\circ\ 27'\ 2''$ & $1^\circ\ 47'\ 48''$ \\
\end{tabular}\endgroup
\end{center}

\par\indent The position of \emph{Juno} calculated for that time from the fourth set of elements was as follows.

\begin{center}
\begingroup\fontsize{9}{11}\selectfont\setlength{\tabcolsep}{3pt}\renewcommand{\arraystretch}{1.1}
\begin{tabular}{c|c|c|c}
& & \multicolumn{2}{c}{Difference}\\ R.A. $\JunoSymbol$ & Decl. bor. $\JunoSymbol$ & in R.A. & in decl.\\ \hline
$30^\circ\ 27'\ 41''6$ & $1^\circ\ 47'\ 20''$ & $+39''6$ & $-27''7$ \\
\end{tabular}\endgroup
\end{center}
\SourceLeaf{270}
\par\indent From that fifth set of elements Dr. \textsc{Gauss} calculated the following ephemeris for the geocentric path of Juno [from 20 October 1805 to 23 July 1806, as for Ceres and Pallas on pp. \hyperlink{native-267}{262} and \hyperlink{native-268}{263} of this volume].

[In the rest of this article, von \textsc{Zach} also gives observations by \textsc{Oriani} in Milan and by \textsc{Olbers} in Bremen, together with a comparison of these observations with \textsc{Gauss}’s fourth set of elements. He also gives the elements of Juno’s orbit calculated by Dr. \textsc{Burckhardt}, from the \emph{Moniteur} of 28 December 1804.]

[\textsc{Bode} published the elements calculated by \textsc{Gauss} and referred to Berlin time, the X\textsuperscript{th} for Ceres, the VIII\textsuperscript{th} for Pallas and the V\textsuperscript{th} for Juno, in the \emph{Astronomisches Jahrbuch für 1808}, Berlin, 1805, \href{https://www.digitale-sammlungen.de/en/view/bsb10538337?page=280}{p. 270}.]

